\documentclass[11pt,a4paper]{article}

\usepackage[utf8]{inputenc}
\usepackage{amsmath,amssymb,amsfonts}
\usepackage{booktabs,tabularx,longtable}
\usepackage{enumitem}
\usepackage{geometry}
\usepackage{hyperref}
\usepackage{xcolor}
\usepackage{listings}
\usepackage{fancyhdr}
\usepackage{titlesec}

\geometry{margin=2.3cm,top=2.8cm,bottom=2.8cm}

\hypersetup{
  colorlinks=true,
  linkcolor=blue!60!black,
  citecolor=blue!60!black,
  urlcolor=blue!60!black,
  pdftitle={ISLS v2.0 Unified Specification},
  pdfauthor={Sebastian Klemm},
}

\lstset{
  basicstyle=\ttfamily\footnotesize,
  breaklines=true,
  frame=single,
  backgroundcolor=\color{gray!8},
  keywordstyle=\color{blue!70!black}\bfseries,
  stringstyle=\color{red!60!black},
  commentstyle=\color{green!50!black}\itshape,
  numbers=left,
  numberstyle=\tiny\color{gray},
  numbersep=6pt,
  tabsize=4,
  showstringspaces=false,
}

\titleformat{\section}{\Large\bfseries}{\thesection}{1em}{}
\titleformat{\subsection}{\large\bfseries}{\thesubsection}{1em}{}
\titleformat{\subsubsection}{\normalsize\bfseries}{\thesubsubsection}{1em}{}

\pagestyle{fancy}
\fancyhf{}
\fancyhead[L]{\small ISLS v2.0 --- Unified Specification}
\fancyhead[R]{\small Klemm, 2026}
\fancyfoot[C]{\thepage}

\begin{document}

\begin{titlepage}
\begin{center}
\vspace*{1.5cm}
{\Huge\bfseries ISLS v2.0}\\[0.4cm]
{\LARGE\bfseries Hypercube Decomposer}\\[0.3cm]
{\LARGE\bfseries with Production-Quality Templates}\\[1.2cm]
{\Large Unified Specification v1.0}\\[1.5cm]
{\large Sebastian Klemm}\\[0.3cm]
{\small March 2026}\\[2cm]

\begin{tabular}{ll}
\textbf{Repository:} & \texttt{graphity/isls} \\
\textbf{New crates:} & \texttt{isls-hypercube}, \texttt{isls-decomposer} \\
\textbf{Modified:} & \texttt{isls-orchestrator} (template upgrade) \\
\textbf{v1.0 preserved:} & \texttt{forge-fullstack} command unchanged \\
\textbf{New command:} & \texttt{forge-v2} \\
\end{tabular}

\vfill
{\small
This specification combines two features into one implementation:\\[0.2cm]
\textbf{Part A}: Rich domain knowledge and production-quality code templates\\
(real entities, JWT auth, pagination, error handling, business rules)\\[0.2cm]
\textbf{Part B}: Recursive hypercube decomposition via Fiedler/Kuramoto/TRITON\\
(replaces fixed 11-stage pipeline with adaptive dimensional collapse)\\[0.4cm]
Both parts are interdependent: the decomposer needs rich templates\\
to emit production code, and the templates need the decomposer\\
to know WHEN and in WHAT ORDER to apply them.
}
\end{center}
\end{titlepage}

\tableofcontents
\newpage

% ============================================================
\part{Architecture}
% ============================================================

\section{Overview}

ISLS v1.0 has a fixed 11-stage pipeline that produces thin CRUD code.
ISLS v2.0 replaces this with:

\begin{enumerate}
\item A \textbf{HyperCube} representation of the requirement space,
  where each dimension is a degree of freedom.
\item A \textbf{Recursive Decomposer} that collapses dimensions using
  Fiedler bisection, Kuramoto synchronization, and TRITON singularity
  detection to determine optimal collapse order.
\item \textbf{Production-quality templates} with domain knowledge that
  emit rich, validated, properly-architected code at each leaf node.
\item A \textbf{learning loop} that crystallizes successful decompositions
  as blueprints for future runs.
\end{enumerate}

\subsection{New Crates}

\begin{center}
\small
\begin{tabular}{lp{9cm}}
\toprule
\textbf{Crate} & \textbf{Responsibility} \\
\midrule
\texttt{isls-hypercube} & HyperCube type, Dimension, Coupling, DimensionValue.
  TOML $\to$ HyperCube parser. Entity inference. Domain knowledge registry
  (warehouse, ecommerce, project-management). Field-level entity templates
  with validations, relationships, indices, business rules. \\
\texttt{isls-decomposer} & Recursive decomposition algorithm. Collapse planner
  (Fiedler bisection, Kuramoto groups, TRITON singularities). Artifact emission
  by calling isls-orchestrator templates. Post-collapse verification via
  isls-reader. Crystallization of successful collapses. CLI entry point
  \texttt{forge-v2}. \\
\bottomrule
\end{tabular}
\end{center}

\texttt{isls-orchestrator} is \textbf{modified}: its Tera templates are
upgraded from thin CRUD to production-quality. The orchestrator becomes the
code-emission backend that the decomposer calls at leaf nodes.

All other existing crates are unchanged. The v1.0 \texttt{forge-fullstack}
command continues to work.

\subsection{Data Flow}

\begin{lstlisting}[title=End-to-End Flow]
warehouse.toml
  | isls-hypercube: parse TOML -> HyperCube (83 dims, 47 free)
  | isls-hypercube: domain enrichment (warehouse template -> rich entities)
  v
HyperCube
  | isls-decomposer: plan_collapse()
  |   Fiedler -> bisect into data/interface clusters
  |   Kuramoto -> group coupled dims (model+validation, api+pagination)
  |   TRITON -> find singularities (framework choice)
  v
CollapsePlan [Singularity, Bisect, Group, Leaf, ...]
  | isls-decomposer: execute_plan()
  |   For each step:
  |     Leaf/Group -> isls-orchestrator templates -> Artifact (code file)
  |     Bisect -> recurse left, recurse right
  |     Singularity -> fix dimension, recompute graph, recurse
  |   After each artifact: isls-reader verifies topology
  v
Vec<Artifact> (67+ files, 5000+ LOC)
  | isls-decomposer: crystallize()
  v
Updated BlueprintRegistry (next run is faster)
\end{lstlisting}

% ============================================================
\part{Part A --- Domain Knowledge and Templates}
% ============================================================

\section{Domain Knowledge System}

\subsection{Architecture}

Domain knowledge lives in \texttt{isls-hypercube} as a registry of
domain templates. Each template maps to an application domain and
contains complete entity definitions, relationships, business rules,
and API features.

\begin{lstlisting}[language=C,title=Domain Registry]
/// Registry of domain knowledge templates.
pub struct DomainRegistry {
    domains: HashMap<String, DomainTemplate>,
}

impl DomainRegistry {
    /// Built-in domains. Loaded at startup.
    pub fn default() -> Self {
        let mut r = Self { domains: HashMap::new() };
        r.register("warehouse", warehouse_domain());
        r.register("ecommerce", ecommerce_domain());
        r.register("project_management", project_management_domain());
        r
    }

    /// Match keywords from app description to a domain.
    /// Returns the best-matching domain template, or None.
    pub fn match_domain(&self, description: &str, module_names: &[String]) -> Option<&DomainTemplate> {
        // Score each domain by keyword overlap
        // "warehouse, inventory, stock, product, order" -> warehouse
        // "shop, cart, checkout, payment" -> ecommerce
        // "project, task, sprint, milestone" -> project_management
    }
}
\end{lstlisting}

\subsection{DomainTemplate Structure}

\begin{lstlisting}[language=C]
pub struct DomainTemplate {
    pub name: String,
    pub keywords: Vec<String>,  // trigger keywords for matching
    pub entities: Vec<EntityTemplate>,
    pub relationships: Vec<Relationship>,
    pub business_rules: Vec<BusinessRule>,
    pub api_features: ApiFeatures,
}

pub struct EntityTemplate {
    pub name: String,
    pub fields: Vec<FieldDef>,
    pub validations: Vec<ValidationRule>,
    pub indices: Vec<IndexDef>,
    pub description: String,
}

pub struct FieldDef {
    pub name: String,
    pub rust_type: String,       // "String", "i64", "Option<String>", "bool"
    pub sql_type: String,        // "VARCHAR(255) NOT NULL", "BIGINT", etc.
    pub nullable: bool,
    pub default_value: Option<String>,  // SQL default
    pub description: String,
}

pub struct ValidationRule {
    pub name: String,
    pub condition: String,       // Rust boolean expression on &self
    pub message: String,         // Error message if validation fails
}

pub struct IndexDef {
    pub name: String,
    pub columns: String,
    pub unique: bool,
}

pub struct Relationship {
    pub from_entity: String,
    pub to_entity: String,
    pub kind: RelKind,
    pub foreign_key: String,
    pub on_delete: String,       // "CASCADE", "SET NULL", "RESTRICT"
}

pub enum RelKind { BelongsTo, HasMany, ManyToMany }

pub struct BusinessRule {
    pub name: String,
    pub description: String,
    pub trigger: String,         // "on_create", "on_stock_change", etc.
    pub logic_pseudocode: String,// Human-readable logic
    pub service_method: String,  // Which service method implements this
    pub entities_involved: Vec<String>,
}

pub struct ApiFeatures {
    pub pagination: bool,
    pub filtering: Vec<String>,  // filterable field names
    pub sorting: Vec<String>,    // sortable field names
    pub search_fields: Option<String>, // comma-separated searchable fields
    pub export_formats: Vec<String>,   // "csv", "json"
}
\end{lstlisting}

\subsection{Warehouse Domain Template}

This is the primary domain template. It must be complete and realistic---
every field, validation, and business rule that a real warehouse system
would have.

\subsubsection{Entities}

\textbf{Product:} id, sku (VARCHAR 50 UNIQUE), name, description (nullable),
category, unit\_price\_cents (BIGINT), cost\_price\_cents, quantity\_on\_hand
(INTEGER DEFAULT 0), reorder\_level (INTEGER DEFAULT 10), reorder\_quantity,
weight\_grams (nullable), barcode (nullable UNIQUE), is\_active (BOOLEAN
DEFAULT true), supplier\_id (FK nullable), warehouse\_id (FK), created\_at,
updated\_at. Validations: sku not empty, name not empty, prices non-negative,
quantity non-negative, reorder level positive.

\textbf{Warehouse:} id, name, code (VARCHAR 20 UNIQUE), address, city,
country, capacity\_units, is\_active, manager\_name (nullable),
contact\_email (nullable), created\_at, updated\_at.

\textbf{Supplier:} id, name, contact\_email, contact\_phone (nullable),
address (nullable), lead\_time\_days (DEFAULT 7), is\_active, rating
(nullable, CHECK 1--5), created\_at, updated\_at.

\textbf{Order:} id, order\_number (VARCHAR 50 UNIQUE), status (VARCHAR 30
DEFAULT `pending'), order\_type (`inbound'/`outbound'/`transfer'),
customer\_name (nullable), customer\_email (nullable), warehouse\_id (FK),
total\_amount\_cents, notes (nullable), shipped\_at (nullable),
delivered\_at (nullable), created\_at, updated\_at. Validations: valid
status enum, valid order type enum.

\textbf{OrderItem:} id, order\_id (FK CASCADE), product\_id (FK),
quantity (INTEGER), unit\_price\_cents (BIGINT), created\_at.
Validation: quantity positive.

\textbf{StockMovement:} id, product\_id (FK), warehouse\_id (FK),
movement\_type (`receipt'/`shipment'/`adjustment'/`transfer'/`return'),
quantity (INTEGER), reference\_type (nullable), reference\_id (nullable),
notes (nullable), created\_at. Validations: valid movement type,
nonzero quantity.

\textbf{User:} id, email (UNIQUE), password\_hash, name, role (DEFAULT
`operator', valid: admin/manager/operator/viewer), is\_active, last\_login\_at
(nullable), created\_at, updated\_at.

\subsubsection{Relationships}

Product $\to$ Warehouse (BelongsTo, warehouse\_id).
Product $\to$ Supplier (BelongsTo, supplier\_id, SET NULL).
Order $\to$ Warehouse (BelongsTo, warehouse\_id).
OrderItem $\to$ Order (BelongsTo, order\_id, CASCADE).
OrderItem $\to$ Product (BelongsTo, product\_id).
StockMovement $\to$ Product (BelongsTo, product\_id).
StockMovement $\to$ Warehouse (BelongsTo, warehouse\_id).

\subsubsection{Business Rules}

\begin{enumerate}
\item \textbf{auto\_reorder}: When quantity\_on\_hand drops below
  reorder\_level and product is active, log a reorder warning.
  Service: \texttt{inventory\_service::check\_reorder\_level()}.

\item \textbf{stock\_on\_fulfillment}: When order status changes to
  `shipped', create outbound StockMovements for each OrderItem and
  decrease Product.quantity\_on\_hand. Service:
  \texttt{order\_service::fulfill\_order()}.

\item \textbf{prevent\_negative\_stock}: Outbound movements cannot
  reduce stock below zero. Return error if insufficient.
  Service: \texttt{inventory\_service::adjust\_stock()}.

\item \textbf{order\_total\_calculation}: When OrderItems change,
  recalculate Order.total\_amount\_cents as sum of item quantities
  times unit prices. Service: \texttt{order\_service::recalculate\_total()}.

\item \textbf{order\_state\_machine}: Orders follow:
  pending $\to$ confirmed $\to$ processing $\to$ shipped $\to$ delivered.
  Also: any $\to$ cancelled (except delivered). Invalid transitions
  return error. Service: \texttt{order\_service::update\_status()}.
\end{enumerate}

\subsubsection{API Features}

Pagination: yes (page, per\_page, default 25, max 100).
Filtering: status, category, warehouse\_id, is\_active, order\_type, movement\_type.
Sorting: created\_at, name, quantity\_on\_hand, unit\_price\_cents, order\_number.
Search: name, sku, description (ILIKE).
Export: CSV, JSON.

\subsection{E-Commerce Domain Template (Abbreviated)}

Entities: Product (name, slug, price, compare\_price, images, inventory\_count,
category\_id), Category (name, slug, parent\_id, position), Customer (email,
name, addresses), Cart (customer\_id, session\_id, items), CartItem (cart\_id,
product\_id, quantity), Order (order\_number, customer\_id, status, subtotal,
tax, shipping, total, payment\_status), OrderLine, Payment (order\_id, provider,
amount, status, transaction\_id), Address, Review (product\_id, customer\_id,
rating, text), Coupon.

Business rules: cart abandonment, inventory reservation on checkout, payment
capture, refund flow, review moderation.

\subsection{Project Management Domain Template (Abbreviated)}

Entities: Project (name, description, owner\_id, status, start\_date, end\_date),
Sprint (project\_id, name, start\_date, end\_date, goal), Task (project\_id,
sprint\_id, title, description, assignee\_id, status, priority, estimate\_hours,
actual\_hours), TeamMember (user\_id, project\_id, role), Comment (task\_id,
author\_id, text), Label, Milestone.

Business rules: sprint velocity calculation, task state machine (todo $\to$
in\_progress $\to$ review $\to$ done), burndown tracking, auto-assignment.

% ============================================================
\section{Production-Quality Code Templates}

These are Tera templates in isls-orchestrator that produce real code, not stubs.
Each template receives context from the domain template (entities, fields,
validations, business rules).

\subsection{Model Template}

For each entity, generate \texttt{src/models/\{entity\_snake\}.rs}:

\begin{itemize}
\item Struct with all fields from EntityTemplate
\item \texttt{\#[derive(Debug, Clone, Serialize, Deserialize, sqlx::FromRow)]}
\item \texttt{Create\{Entity\}} payload struct (no id, no timestamps)
\item \texttt{Update\{Entity\}} payload struct (all fields Optional)
\item \texttt{impl \{Entity\} \{ fn validate(\&self) -> Result<(), Vec<String>> \}}
  with all validation rules
\item \texttt{impl Create\{Entity\} \{ fn validate(\&self) -> Result<(), Vec<String>> \}}
\end{itemize}

\subsection{Migration Template}

Generate \texttt{database/migrations/001\_initial.sql}:
\begin{itemize}
\item CREATE TABLE for each entity with all fields, types, constraints
\item Foreign key constraints with ON DELETE
\item CREATE INDEX for each index definition
\item Ordered by dependency (referenced tables first)
\end{itemize}

\subsection{Query Template}

For each entity, generate \texttt{src/database/\{entity\_snake\}\_queries.rs}:
\begin{itemize}
\item \texttt{list\_all(pool, params: PaginationParams) -> PaginatedResponse}
  with LIMIT/OFFSET, ORDER BY, WHERE filters, ILIKE search
\item \texttt{get\_by\_id(pool, id) -> Result<Entity>}
\item \texttt{insert(pool, payload: CreateEntity) -> Result<Entity>}
  with RETURNING clause
\item \texttt{update(pool, id, payload: UpdateEntity) -> Result<Entity>}
  with COALESCE for partial updates
\item \texttt{delete(pool, id) -> Result<bool>}
\item Additional queries per business rule (e.g., \texttt{find\_by\_sku},
  \texttt{find\_low\_stock}, \texttt{count\_by\_status})
\item All use runtime \texttt{sqlx::query\_as::<\_, Entity>()} --- no
  compile-time macros
\end{itemize}

\subsection{Error Template}

Generate \texttt{src/errors.rs} (one file, shared):
\begin{itemize}
\item \texttt{AppError} enum: NotFound, BadRequest, ValidationError(Vec<String>),
  Unauthorized, Forbidden, DatabaseError, InternalError, Conflict
\item \texttt{impl Display} and \texttt{impl ResponseError} with proper
  HTTP status codes and JSON error bodies
\item \texttt{impl From<sqlx::Error>} with RowNotFound $\to$ NotFound mapping
\end{itemize}

\subsection{Pagination Template}

Generate \texttt{src/pagination.rs} (one file, shared):
\begin{itemize}
\item \texttt{PaginationParams}: page, per\_page, sort, order, search (all Optional)
\item \texttt{PaginatedResponse<T>}: data Vec, pagination meta
\item \texttt{PaginationMeta}: page, per\_page, total, total\_pages
\item Helper methods: offset(), clamped per\_page (1--100)
\end{itemize}

\subsection{Auth Template}

Generate \texttt{src/auth.rs} (one file, shared):
\begin{itemize}
\item \texttt{Claims} struct (sub, email, role, exp, iat)
\item \texttt{create\_token(user, secret) -> Result<String>}
\item \texttt{verify\_token(token, secret) -> Result<Claims>}
\item \texttt{AuthUser} extractor (from Authorization Bearer header)
\item \texttt{require\_role(user, required) -> Result<()>} with hierarchy:
  viewer $<$ operator $<$ manager $<$ admin
\item Login endpoint: verify password with bcrypt, return JWT
\item Register endpoint: hash password, create user
\end{itemize}

Dependencies to add to generated Cargo.toml: \texttt{jsonwebtoken = "9"},
\texttt{bcrypt = "0.15"}.

\subsection{Service Template}

For each entity, generate \texttt{src/services/\{entity\_snake\}.rs}:
\begin{itemize}
\item CRUD methods: create, get, list, update, delete
\item Each method: role check, input validation, business rule enforcement,
  query call, error mapping
\item Additional methods per business rule (adjust\_stock, fulfill\_order,
  update\_status with state machine validation)
\item All return \texttt{Result<T, AppError>}
\end{itemize}

\subsection{API Template}

For each entity, generate \texttt{src/api/\{entity\_snake\}.rs}:
\begin{itemize}
\item \texttt{configure(cfg)} registers all routes
\item GET \texttt{/api/\{entities\}} $\to$ list (paginated, with query params)
\item GET \texttt{/api/\{entities\}/\{id\}} $\to$ get by id
\item POST \texttt{/api/\{entities\}} $\to$ create (requires auth)
\item PUT \texttt{/api/\{entities\}/\{id\}} $\to$ update (requires auth)
\item DELETE \texttt{/api/\{entities\}/\{id\}} $\to$ delete (requires admin/manager)
\item Additional routes per business rule (POST \texttt{/api/orders/\{id\}/fulfill},
  POST \texttt{/api/products/\{id\}/adjust-stock})
\item Auth endpoints: POST \texttt{/api/auth/login}, POST \texttt{/api/auth/register}
\item All handlers: extract AuthUser, call service, return JSON with proper status codes
\end{itemize}

\subsection{Config Template}

Generate \texttt{src/config.rs}:
\begin{itemize}
\item \texttt{AppConfig}: database\_url, jwt\_secret, server\_host, server\_port,
  log\_level
\item Load from environment variables with defaults
\end{itemize}

\subsection{Main Template}

Generate \texttt{src/main.rs}:
\begin{itemize}
\item Initialize tracing/logging
\item Load config from environment
\item Create database connection pool
\item Build Actix-web app with CORS, all route configurations wired
\item Bind and run server
\end{itemize}

\subsection{Frontend Templates}

\subsubsection{Layout and Navigation}

\texttt{frontend/index.html}: HTML shell with nav, content area, CSS link, JS loader.
\texttt{frontend/style.css}: Clean grid/flexbox layout, responsive, table styles,
form styles, button styles, status badges, KPI cards, chart containers.
Color palette: professional blues/grays, accent for actions.

\subsubsection{API Client}

\texttt{frontend/src/api/client.js}: fetch wrapper with JWT header injection,
error handling, base URL configuration, JSON parsing.

\subsubsection{Pages}

For each entity module, generate \texttt{frontend/src/pages/\{module\}.js}:
\begin{itemize}
\item Paginated table with column headers, sorting arrows, search input
\item Filter dropdowns (status, category, etc.)
\item Create/Edit modal or form with field validation
\item Delete confirmation dialog
\item Status badges for status fields
\item Loading spinners, error messages
\end{itemize}

\texttt{frontend/src/pages/dashboard.js}:
\begin{itemize}
\item KPI cards: total products, total orders, revenue, low stock alerts
\item Charts: orders per day (bar), revenue trend (line), stock by category (pie)
\item Recent orders table (last 10)
\item Low stock alerts list
\end{itemize}

\texttt{frontend/src/pages/login.js}: email/password form, JWT storage in
localStorage, redirect to dashboard on success.

\subsubsection{Components}

\texttt{frontend/src/components/table.js}: reusable sortable table.
\texttt{frontend/src/components/form.js}: reusable validated form.
\texttt{frontend/src/components/chart.js}: simple canvas-based charts
(bar, line, pie --- no external library, keep it self-contained).
\texttt{frontend/src/components/navbar.js}: navigation with role-based
menu visibility.
\texttt{frontend/src/components/modal.js}: reusable dialog.
\texttt{frontend/src/components/pagination.js}: page controls.

\subsection{Test Template}

Generate \texttt{backend/tests/api\_tests.rs}:
\begin{itemize}
\item Test helpers: create test app, create test user, get auth token
\item For each entity: test create (201), test list with pagination,
  test get by id (200), test get nonexistent (404), test update,
  test delete, test unauthorized access (401)
\item For each business rule: test the rule works (e.g., stock adjustment
  prevents negative), test invalid state transitions
\end{itemize}

Note: tests use \texttt{sqlx::test} if available, otherwise standard
\texttt{\#[test]} with mock data.

\subsection{Deployment Templates}

\texttt{Dockerfile}: multi-stage build, cargo build release, slim runtime.
\texttt{docker-compose.yml}: postgres + backend + frontend (static files via nginx).
\texttt{.env.example}: all config variables with comments.
\texttt{README.md}: project description, setup instructions, API documentation
with all endpoints listed, example curl commands.

% ============================================================
\part{Part B --- Hypercube Decomposer}
% ============================================================

\section{HyperCube Type System}

\subsection{Core Types}

\begin{lstlisting}[language=C]
pub struct HyperCube {
    pub dimensions: Vec<Dimension>,
    pub couplings: Vec<Coupling>,
    pub depth: u32,
    pub parent_signature: Option<String>,
}

pub struct Dimension {
    pub name: String,           // "model.inventory.Product"
    pub category: DimCategory,
    pub state: DimState,
    pub complexity: u32,        // estimated LOC this dimension produces
    pub description: String,
}

pub enum DimCategory {
    Architecture,   // framework, runtime, monolith/micro (collapse first)
    DataModel,      // entities, fields, types (collapse second)
    Storage,        // database, queries, migrations
    BusinessLogic,  // domain rules, state machines, workflows
    Interface,      // API endpoints, protocols
    Security,       // auth, encryption, roles
    Presentation,   // frontend pages, components, styling
    Testing,        // test strategies, test cases
    Deployment,     // docker, configs, CI/CD
    Documentation,  // README, API docs
}

pub enum DimState {
    Fixed(DimValue),
    Free { options: Option<Vec<DimValue>>, default: Option<DimValue> },
    Derived { depends_on: String },
}

pub enum DimValue {
    Choice(String),
    Code(String),
    EntityDef(EntityTemplate),
    Composite(Vec<DimValue>),
}

pub struct Coupling {
    pub from: String,
    pub to: String,
    pub strength: f64,      // 0.0 to 1.0
    pub direction: CouplingDir,
}

pub enum CouplingDir { Forward, Backward, Mutual }

impl HyperCube {
    pub fn dof(&self) -> usize { /* count Free dims */ }
    pub fn coupling_graph(&self) -> CouplingGraph { /* build adjacency */ }
    pub fn fix(&mut self, name: &str, value: DimValue);
    pub fn extract_subcube(&self, dims: &[String]) -> HyperCube;
    pub fn signature(&self) -> String { /* SHA-256 of canonical state */ }
    pub fn free_dimensions(&self) -> Vec<&Dimension>;
}
\end{lstlisting}

\subsection{CouplingGraph}

\begin{lstlisting}[language=C]
/// The coupling graph over free dimensions.
/// Used for Fiedler bisection and Kuramoto grouping.
pub struct CouplingGraph {
    pub nodes: Vec<String>,            // dimension names
    pub edges: Vec<(usize, usize, f64)>, // (from, to, weight)
    pub adjacency: Vec<Vec<(usize, f64)>>,
}

impl CouplingGraph {
    /// Compute the Laplacian matrix L = D - A
    pub fn laplacian(&self) -> Vec<Vec<f64>>;

    /// Compute Fiedler value (lambda_2) and Fiedler vector
    pub fn fiedler(&self) -> FiedlerResult;

    /// Bisect using Fiedler vector (positive/negative partition)
    pub fn fiedler_bisect(&self) -> (Vec<String>, Vec<String>);

    /// Find Kuramoto-synchronized groups
    /// Dimensions with phase difference < threshold form a group
    pub fn kuramoto_groups(&self, kappa: f64, threshold: f64) -> Vec<DimGroup>;

    /// Detect singularities: dimensions where fixing them
    /// causes the largest spectral gap change
    pub fn singularities(&self, gap_threshold: f64) -> Vec<Singularity>;
}

pub struct FiedlerResult {
    pub value: f64,
    pub vector: Vec<f64>,
}

pub struct DimGroup {
    pub dimensions: Vec<String>,
    pub avg_coupling: f64,
    pub estimated_loc: usize,
}

pub struct Singularity {
    pub dimension: String,
    pub impact: f64,         // fraction of graph affected by fixing this dim
    pub options: Vec<DimValue>,
}
\end{lstlisting}

Note on implementation: The Fiedler computation requires eigenvalue
decomposition. For small graphs ($< 200$ nodes), use power iteration
or QR algorithm implemented directly. Do NOT add LAPACK or nalgebra
as dependencies --- implement a simple eigensolver for the second-smallest
eigenvalue. The existing code in \texttt{isls-topology} may have a Fiedler
computation that can be reused.

For Kuramoto: simulate coupled oscillators on the graph. Each dimension
is an oscillator. Natural frequencies proportional to complexity.
Coupling through graph edges. After convergence, dimensions with similar
phase form groups. The existing code in \texttt{isls-scale} may have
Kuramoto that can be reused.

For singularities: for each free dimension, temporarily remove it from
the graph and measure the change in spectral gap. Dimensions whose
removal causes the largest gap change are singularities.

\section{TOML to HyperCube Parser}

\begin{lstlisting}[language=C]
/// Parse a TOML file into a HyperCube.
pub fn parse_toml_to_cube(toml_str: &str) -> Result<HyperCube> {
    // 1. Parse TOML
    // 2. Extract [app] -> Architecture dims (fixed)
    // 3. Extract [app.modules] -> for each module:
    //      a. Infer entities from description (nouns)
    //      b. Match domain (warehouse/ecommerce/pm) from keywords
    //      c. Enrich entities with domain template fields
    //      d. Create dims: model.X, validation.X, queries.X,
    //         service.X, api.X, tests.X for each entity
    //      e. Infer business rules -> business_logic dims
    // 4. Extract [backend] -> Runtime dims (fixed)
    // 5. Extract [frontend] -> Presentation dims
    // 6. Extract [deployment] -> Deployment dims
    // 7. Add cross-cutting dims: error_handling, pagination, auth, config
    // 8. Auto-generate couplings:
    //      model -> validation (0.95, Forward)
    //      model -> queries (0.90, Forward)
    //      queries -> service (0.85, Forward)
    //      service -> api (0.80, Forward)
    //      model -> tests (0.70, Forward)
    //      auth -> api (0.75, Mutual)
    //      error_handling -> service (0.70, Mutual)
    //      pagination -> api (0.85, Mutual)
    //      config -> main (0.90, Forward)
    // 9. Return HyperCube
}
\end{lstlisting}

\section{Recursive Decomposition}

\begin{lstlisting}[language=C,title=The Core Algorithm]
pub fn decompose(
    cube: &mut HyperCube,
    ctx: &mut DecompositionContext,
) -> Result<Vec<Artifact>> {
    // Base: fully constrained
    if cube.dof() == 0 {
        return Ok(vec![]);
    }

    // Base: matches a known blueprint
    if let Some(bp) = ctx.blueprints.match_cube(cube, 0.80) {
        return emit_from_blueprint(bp, cube, ctx);
    }

    // Base: leaf dimension (single free dim, low complexity)
    if cube.dof() == 1 {
        return emit_leaf(cube, ctx);
    }

    // Recursive: plan and execute collapse
    let graph = cube.coupling_graph();
    let plan = plan_collapse(&graph, &ctx.blueprints);

    let mut artifacts = Vec::new();

    for step in plan.steps {
        match step {
            CollapseStep::Singularity { dimension, .. } => {
                // Fix the architectural decision
                let value = resolve_singularity(&dimension, cube, ctx)?;
                cube.fix(&dimension, value);
                // Graph changed fundamentally -- restart decomposition
                let sub = decompose(cube, ctx)?;
                artifacts.extend(sub);
                return Ok(artifacts);
            }
            CollapseStep::Group { dimensions, .. } => {
                // Extract sub-cube, decompose, fix dims with result
                let mut sub_cube = cube.extract_subcube(&dimensions);
                let sub_artifacts = emit_group(&mut sub_cube, ctx)?;
                for dim in &dimensions {
                    cube.fix(dim, DimValue::Choice("generated".into()));
                }
                // Verify topology of generated code
                verify_artifacts(&sub_artifacts, ctx)?;
                artifacts.extend(sub_artifacts);
            }
            CollapseStep::Bisect { left, right, .. } => {
                let mut left_cube = cube.extract_subcube(&left);
                let mut right_cube = cube.extract_subcube(&right);
                artifacts.extend(decompose(&mut left_cube, ctx)?);
                artifacts.extend(decompose(&mut right_cube, ctx)?);
                for dim in left.iter().chain(right.iter()) {
                    cube.fix(dim, DimValue::Choice("generated".into()));
                }
            }
            CollapseStep::Leaf { dimension, method } => {
                let art = emit_leaf_dim(&dimension, &method, cube, ctx)?;
                cube.fix(&dimension, DimValue::Choice("generated".into()));
                artifacts.extend(art);
            }
        }
    }

    // Crystallize this decomposition for future reuse
    let crystal = crystallize(cube, &artifacts);
    ctx.blueprints.add(crystal);

    Ok(artifacts)
}
\end{lstlisting}

\subsection{Collapse Planner}

\begin{lstlisting}[language=C]
pub fn plan_collapse(
    graph: &CouplingGraph,
    blueprints: &BlueprintRegistry,
) -> CollapsePlan {
    let mut steps = Vec::new();

    // 1. Singularities first
    let sings = graph.singularities(0.01);
    for s in sings {
        steps.push(CollapseStep::Singularity {
            dimension: s.dimension,
            options: s.options,
            impact: s.impact,
        });
    }
    // If singularities found, return early -- graph will be rebuilt
    if !steps.is_empty() {
        return CollapsePlan { steps, .. };
    }

    // 2. Try Fiedler bisection
    let fiedler = graph.fiedler();
    if graph.nodes.len() > 3 && fiedler.value < 0.5 {
        // Clean bisection possible
        let (left, right) = graph.fiedler_bisect();
        if !left.is_empty() && !right.is_empty() {
            steps.push(CollapseStep::Bisect {
                left, right, fiedler_value: fiedler.value,
            });
            return CollapsePlan { steps, .. };
        }
    }

    // 3. Kuramoto groups for tightly coupled graphs
    let groups = graph.kuramoto_groups(2.0, 0.8);
    for group in groups {
        if group.dimensions.len() > 1 {
            steps.push(CollapseStep::Group {
                dimensions: group.dimensions,
                kuramoto_phase: group.avg_coupling,
                estimated_loc: group.estimated_loc,
            });
        }
    }

    // 4. Remaining isolated dims as leaves
    let grouped: HashSet<_> = steps.iter()
        .flat_map(|s| s.dimensions())
        .collect();
    for dim in graph.nodes.iter() {
        if !grouped.contains(dim) {
            let method = if let Some(bp) = blueprints.match_dim(dim) {
                ResolutionMethod::Blueprint(bp.id.clone())
            } else {
                ResolutionMethod::Template(dim.clone())
            };
            steps.push(CollapseStep::Leaf {
                dimension: dim.clone(), method,
            });
        }
    }

    CollapsePlan { steps, .. }
}
\end{lstlisting}

\subsection{Emission: Connecting to Templates}

When the decomposer reaches a leaf or group, it calls isls-orchestrator's
template system:

\begin{lstlisting}[language=C]
/// Emit code for a group of related dimensions.
fn emit_group(
    sub_cube: &mut HyperCube,
    ctx: &mut DecompositionContext,
) -> Result<Vec<Artifact>> {
    let mut artifacts = Vec::new();

    // Determine what kind of group this is based on dimension categories
    let categories: Vec<_> = sub_cube.free_dimensions()
        .iter().map(|d| &d.category).collect();

    if categories.iter().all(|c| matches!(c, DimCategory::DataModel)) {
        // Model group: emit model files from domain template
        for dim in sub_cube.free_dimensions() {
            if let DimState::Free { .. } = &dim.state {
                let entity = ctx.domain.find_entity(&dim.name)?;
                let code = ctx.templates.render_model(&entity)?;
                artifacts.push(Artifact::SourceFile {
                    path: format!("backend/src/models/{}.rs", entity.name.to_snake()),
                    language: "rust".into(),
                    content: code,
                    evidence_hash: sha256(&code),
                });
            }
        }
    } else if categories.iter().all(|c| matches!(c, DimCategory::Storage)) {
        // Database group: emit queries + migrations
        // ... similar pattern
    } else if categories.iter().all(|c| matches!(c, DimCategory::Interface)) {
        // API group: emit endpoints with pagination, auth
        // ... similar pattern
    } else if categories.iter().all(|c| matches!(c, DimCategory::Presentation)) {
        // Frontend group: emit pages, components
        // ... similar pattern
    } else {
        // Mixed group: emit each dimension individually
        for dim in sub_cube.free_dimensions() {
            artifacts.extend(emit_leaf_dim(&dim.name, &ResolutionMethod::Template(dim.name.clone()), sub_cube, ctx)?);
        }
    }

    Ok(artifacts)
}
\end{lstlisting}

% ============================================================
\part{Integration and Verification}
% ============================================================

\section{CLI Entry Point}

\begin{lstlisting}[language=bash]
# New v2.0 command
isls forge-v2 --requirements examples/warehouse.toml \
              --output ./output/warehouse \
              --mock-oracle \
              --trace            # print decomposition trace

# v1.0 still works
isls forge-fullstack --constraints examples/warehouse.toml --output ./out
\end{lstlisting}

The \texttt{--trace} flag prints the full decomposition tree: which
dimensions were collapsed, in what order, by which method (bisect/group/
leaf), which blueprints were hit, and the final artifact list.

\section{Decomposition Trace Output}

\begin{lstlisting}[title=Expected Output with --trace]
=== ISLS v2.0 Hypercube Decomposer ===

Parsing warehouse.toml...
  Domain detected: warehouse (keywords: inventory, stock, order)
  Entities enriched: Product(17 fields), Warehouse(12), Supplier(10),
    Order(13), OrderItem(6), StockMovement(9), User(9)
  Business rules: 5

HyperCube: 83 dimensions (36 fixed, 47 free)
Couplings: 156 edges

=== Decomposition ===

[depth=0] Singularity: runtime.framework
  Impact: 34% | Fixed: "actix-web"
  Rebuilding coupling graph...

[depth=0] Fiedler bisect (lambda_2=0.23)
  Left (22 dims): model.*, validation.*, queries.*, migration.*
  Right (25 dims): service.*, api.*, auth.*, frontend.*, deployment.*

  [depth=1, LEFT] Kuramoto group (phase=0.91)
    Dims: model.inventory.Product, validation.inventory.Product, ...
    Emitting: 7 model files, 7 validation modules
    Method: domain template (warehouse)
    LOC: 1,240

  [depth=1, LEFT] Kuramoto group (phase=0.87)
    Dims: queries.*, migration.*
    Emitting: 7 query files, 1 migration
    Method: template (paginated queries)
    LOC: 890

  [depth=1, RIGHT] Fiedler bisect (lambda_2=0.31)
    Left: service.*, api.*, auth.*, error.*, pagination.*
    Right: frontend.*, deployment.*, docs.*

    [depth=2, LEFT] Kuramoto group (phase=0.88)
      Dims: error_handling, pagination, auth, config
      Emitting: errors.rs, pagination.rs, auth.rs, config.rs
      Method: template (production-quality)
      LOC: 680

    [depth=2, LEFT] Group: service.*, api.*
      Emitting: 7 service files, 7 API files
      Method: template (with business rules injected)
      LOC: 1,450

    [depth=2, RIGHT] Group: frontend.*
      Emitting: 5 pages, 6 components, API client, styles
      Method: template
      LOC: 1,180

    [depth=2, RIGHT] Leaf: deployment.*
      Emitting: Dockerfile, docker-compose.yml, .env, README
      Method: template
      LOC: 280

  [depth=1] Leaf: testing.*
    Emitting: api_tests.rs, service_tests.rs
    Method: template (per business rule)
    LOC: 520

=== Verification ===
  Barbara topology check: 67 files parsed
  Average similarity to plan: 0.84
  All files > 0.70: YES

=== Crystallization ===
  New blueprints: 8
  Updated blueprints: 31
  Registry total: 39

=== Result ===
  Files:           67
  LOC:             6,240
  Oracle calls:    0 (mock mode)
  Blueprint hits:  31
  Template hits:   36
  Time:            0.8s
  Evidence chain:  VALID (83 entries)
\end{lstlisting}

\section{Verification Pipeline}

After all artifacts are emitted, the decomposer runs verification:

\begin{enumerate}
\item \textbf{Structural verification}: every API endpoint has a
  matching service function. Every service uses existing queries.
  Every model referenced in queries exists. Every frontend API call
  hits an existing endpoint.

\item \textbf{Topological verification} (if isls-reader available):
  parse generated Rust files with tree-sitter, compute topology,
  compare against expected structure from domain template.

\item \textbf{Compilation verification} (if possible): attempt
  \texttt{cargo build} on the generated backend.
\end{enumerate}

Note: compilation verification requires running cargo in a subprocess.
If this fails or is unavailable, skip it and log a warning. The
structural verification must always run.

\section{Success Criteria}

\begin{enumerate}
\item \texttt{cargo build --workspace} --- zero errors, all existing tests pass.
\item \texttt{isls forge-v2 --requirements examples/warehouse.toml --mock-oracle --output ./out --trace}
  runs successfully and prints the decomposition trace.
\item The trace shows Fiedler bisection, Kuramoto groups, and dimensional collapse.
\item Generated backend compiles with \texttt{cargo build} (zero errors).
\item Generated code has rich domain models: Product with $\geq 15$ fields.
\item Generated code has real error handling (AppError with $\geq 6$ variants).
\item Generated code has pagination (PaginationParams, PaginatedResponse).
\item Generated code has JWT auth (create\_token, verify\_token, AuthUser extractor).
\item Generated code has $\geq 3$ business rules implemented in services
  (stock adjustment, order fulfillment, state machine).
\item Generated LOC $\geq 5,000$ (up from 2,462 in v1.0).
\item Generated files $\geq 60$ (up from 48 in v1.0).
\item Second run of same spec is faster (more blueprint hits in trace).
\item v1.0 \texttt{forge-fullstack} still works unchanged.
\end{enumerate}

\section{Implementation Constraints for Claude Code}

\begin{enumerate}
\item \textbf{Repository: graphity/isls.} Add isls-hypercube and
  isls-decomposer to workspace members.

\item \textbf{Reuse existing PSE crates.} Read the actual APIs of
  isls-topology (Fiedler), isls-scale (Kuramoto), isls-navigator
  (TRITON/singularity). If their APIs don't match the pseudocode
  in this spec, implement minimal local versions in
  isls-decomposer. Do NOT modify existing crates unless necessary.

\item \textbf{Template upgrade goes in isls-orchestrator.} The
  production-quality templates (Part A) replace the existing thin
  templates. Keep both: a ``v1'' template set for \texttt{forge-fullstack}
  backward compatibility, and a ``v2'' template set for \texttt{forge-v2}.
  Or upgrade in-place if the v2 templates are strictly better.

\item \textbf{Domain templates are data.} The warehouse/ecommerce/pm
  domain knowledge is defined as Rust structs in isls-hypercube,
  not as Tera templates. The Tera templates receive domain data
  as context.

\item \textbf{Sequential implementation:}
  \begin{enumerate}
  \item isls-hypercube: types + TOML parser + domain registry + warehouse domain.
    Test: parse warehouse.toml into HyperCube with 83 dims.
  \item isls-decomposer: collapse planner with Fiedler/Kuramoto.
    Test: plan\_collapse produces correct step sequence for warehouse cube.
  \item Template upgrade in isls-orchestrator: production-quality templates.
    Test: render model template for warehouse Product produces 15+ fields.
  \item Wire decomposer to orchestrator: emit\_group calls templates.
    Test: full warehouse decomposition produces 60+ files.
  \item Verification: structural checks post-generation.
  \item Crystallization: blueprints saved after successful run.
  \item CLI: forge-v2 command.
  \item End-to-end: warehouse.toml $\to$ compilable backend with rich code.
  \end{enumerate}

\item \textbf{Mock Oracle.} All generation uses templates. No LLM calls.
  The \texttt{--mock-oracle} flag must produce all output without network.

\item \textbf{All existing tests must pass.} Zero regressions.

\item \textbf{The decomposition trace is the most important deliverable
  after compilable output.} If the trace shows Fiedler bisection and
  Kuramoto groups producing correct code, the algorithm is proven.

\item \textbf{No \texttt{unwrap()} in library code.}
\item \textbf{Every public item documented.}
\item \textbf{MIT License. Copyright (c) 2026 Sebastian Klemm.}
\end{enumerate}

\vfill
\noindent\rule{\textwidth}{0.4pt}
\begin{center}
\small
ISLS v2.0 Unified Specification --- Sebastian Klemm --- March 2026\\
Hypercube Decomposer + Production-Quality Templates\\
Every application is a hypercube. Generation is dimensional collapse.\\
PSE navigates. Barbara verifies. The system learns.\\
\url{https://github.com/LashSesh/graphity}
\end{center}

\end{document}
